% s-210-p210-m002; original PDF 210
\Aphorism{s-210-p210-m002}{303}
„So ist’s! so habe ich’s erlebt!“ ruft der Leser, wenn eine besonders ausgedachte Feinheit sich gerade an ihn zu wenden, gerade auf seinen Fall zu beziehen scheint. Mit dieser persönlichen Quittung glaubt er dem Autor zu schmeicheln und Psychologie anzurühmen. Aber was Jemand erfährt, ist nur ein Theil und eine Entstellung von dem, was er erfahren könnte; Jeder hat seinen besonderen Modus, Erlebnisse zu versäumen und zu verfälschen. Irgend etwas Schiefes und Krummes ist in allem Erleben; hat man es nun glücklich aus dem Gröbsten herausgearbeitet und Sinn, Ordnung, gerade Linie hineingelegt, d. h. hat man es \Sperr{gedichtet}, so ist es ein Rückschritt, die Erinnerung an das rohe Original nicht loswerden zu können und immer wieder das überwundene Persönliche hervorzusuchen. Der Autor hat sein Brutto-Erlebniss stilisirt, gereinigt, durchgefiltert: soll das der Lohn seiner Mühe sein, dass der erste beste Mitmensch von Leser ihm hinter die Coulissen sieht und in dem persönlich missverstandenen Dichter sich getroffen fühlt?
\par

